\subsection{Cyclic auxiliary repair at every finite three-power precision}

\begin{statement}{A.130}{Proved}
{The exact higher-precision local image for the relative norm torus}
Let \(a\in\Q^\times\) satisfy \(v_3(a)=1\), let \(n,r\ge1\),
and use the torus and coefficient maps of A.129.  Put
\[
 \begin{gathered}
 N=3^n,\quad d=N/3,\quad R=\Z/3^r\Z,\quad C=C_r=T_n[3^r],\\
 L=\Q(a^{1/N},\mu_{3^{\max(n,r)}}),\qquad
                      G=\operatorname{Gal}(L/\Q).
 \end{gathered}
 \tag{A.130.1}
\]
Let \(E/\Q\) be finite Galois and contain \(L\), and let
\(S\) be finite and contain every prime ramified in \(E\),
the prime three and the real place.  Define
\[
 \begin{split}
 \mathcal H_r&=\bigcup_{\substack{P\subset{
m cs}(E)\setminus S\\
                                  P\ {\rm finite}}}
                         H^1(G_{\Q,S\cup P},C),\\
 \mathcal H_{r,\rm cyc}&=\{\gamma\in\mathcal H_r:
       \operatorname{res}_p\gamma\text{ has cyclic image at every}
                        \ p\notin S\text{ where it ramifies}\}.
 \end{split}
 \tag{A.130.2}
\]
Every such ramified prime splits the coefficient field, so
the local image is that of an actual homomorphism into \(C\).
Let \(\mathcal H_{r,\rm in}\subset\mathcal H_{r,\rm cyc}\)
be the classes for which the inertia image generates that
local image at every new ramified prime.  In tame coordinates
this is the exact condition \(c_p=z_pt_p\) for some \(z_p\in R\).
For the \(n\) forms \(B_j\) of A.115, let
\[
 Q_{j,w}(\eta)=\operatorname{inv}_w\!\left(
       \tfrac12(B_j)_*((\rho_r)_*\eta\cup(\rho_r)_*\eta)\right).
 \tag{A.130.3}
\]
Then the exact image equality is
\[
 \operatorname{loc}_S\mathcal H_{r,\rm in}
   =\operatorname{loc}_S\mathcal H_{r,\rm cyc}
   =\left\{(\eta_w)\in\operatorname{loc}_S\mathcal H_r:
           \sum_{w\in S}Q_{j,w}(\eta_w)=0
                                  \text{ for every }0\le j<n\right\}.
 \tag{A.130.4}
\]
Thus, at every fixed precision, the additional reciprocity
conditions are exactly the mod-three conditions under the
displayed coefficient map.  All old local classes are
preserved at their original precision \(3^r\).

In particular, arbitrary prescriptions at finitely many
split-\(E\) primes, with zero prescribed at the other
controlled places, have such a lift exactly when their
total values (A.130.3) vanish.  Pure tame prescriptions
always have this property.  When every prescribed class
on \(S\) is an actual local Kummer class for \(T_n\), the
global class may be represented by an actual
\[
                 U\in K_n^\times,\qquad
                       N_{K_n/K_{n-1}}U=1.
 \tag{A.130.5}
\]
The theorem prescribes a finite precision and a finite local
tuple.  It does not assert that the classes chosen for
different precisions reduce to one another.
\end{statement}

\paragraph{The odd-prime-power cyclic-support lemma.}
We use the following extension of the odd-prime argument in
\cite[Theorem~13 and its proof, pp.~14--19]{RootSchmidtWingberg1998}.
Let \(p\) be odd, let \(m=p^r\), and let
\(\Omega/K/k\) be finite extensions of number fields with both
\(\Omega/k\) and \(K/k\) Galois.  Suppose \(\mu_m\subset K\).
Let \(A\) be a finite \(\mathbf Z/m[\operatorname{Gal}(K/k)]\)-module,
and let \(T\) contain the primes ramified in \(\Omega/k\),
the primes above \(p\), and the infinite places.  Put
\[
 S=\operatorname{cs}(\Omega/k)\cup T.
 \tag{A.130.S1}
\]
Here \(k_S\) is the maximal extension unramified outside \(S\);
all cohomology is continuous.  Suppose
\(y\in H^1(k_S/K,A)\) is unramified at every place of \(T(K)\)
and is zero at the places above
\(\operatorname{Ram}(K/k)\cup S_p\cup S_\infty\).
There exists \(z\in H^1(k_S/K,A)\) with the following exact
properties:
\[
\begin{aligned}
 z_{\mathfrak p}&=y_{\mathfrak p}
                         &&(\mathfrak p\in T(K)),\\
 z_{\mathfrak p}\text{ ramified},\ \mathfrak p\notin T(K)
 &\ \Longrightarrow\
 \begin{cases}
 z_{\mathfrak p}\text{ is split by a cyclic extension of }K_{\mathfrak p},\\
 z_{\sigma\mathfrak p}=0
       \quad\text{for all }1\ne\sigma\in\operatorname{Gal}(K/k).
 \end{cases}
\end{aligned}
 \tag{A.130.S2}
\]
Consequently \(x=\operatorname{cor}_{K/k}z\) satisfies
\[
 x_v=(\operatorname{cor}_{K/k}y)_v\quad(v\in T),
 \qquad x_v\text{ is cyclic}\quad(v\notin T).
 \tag{A.130.S3}
\]
The cyclicity assertion means that the actual local cohomology
class becomes zero over one finite cyclic extension.
In particular, at a place splitting \(K\), it means that the
image of its local homomorphism is a cyclic subgroup of \(A\).
The construction gives a stronger assertion when the underlying
group \(A\) is free over \(\mathbf Z/m\): at every new ramified
place the inertia image has exact order \(m\) and generates
the entire local image.  In tame coordinates \(t,c\in A\)
this gives the literal relation \(c=zt\) for some
\(z\in\mathbf Z/m\), with \(t\) of order \(m\).

\emph{The finite duality input.}
We spell out the part of the proof where a change of coefficient
exponent could otherwise be hidden.  For a finite Galois field
\(\Omega_0/K/k\), a set \(T_0\) as above, and
\(S_0=\operatorname{cs}(\Omega_0/k)\cup T_0\), the image of
\[
 H^1(k_{S_0}/K,\mu_e)\longrightarrow
       \prod_{\mathfrak p\in T_0(K)}H^1(K_{\mathfrak p},\mu_e),
                       \qquad e\mid m,
 \tag{A.130.S4}
\]
is the annihilator, under the sum of local Tate pairings, of
\(\operatorname{Hom}(\operatorname{Gal}(\Omega_0/K),\mathbf Z/e)\).
This is the finite-coefficient Poitou--Tate sequence with
unramified conditions outside \(T_0\), followed by Chebotarev.
For clarity, in a finite enlargement of \(T_0\) the annihilating
dual classes are unramified outside \(T_0\), because their
local restrictions at every added place are zero.  They
therefore lie in the one finite group
\(H^1(k_{T_0}/K,\mathbf Z/e)\).
Intersecting their kernels as added primes run through
\(\operatorname{cs}(\Omega_0/k)\) stabilizes in this finite group.
Chebotarev identifies the stabilized subgroup with exactly
the characters factoring through \(\Omega_0/K\): any character
not factoring through that field is nonzero on some Frobenius
whose underlying \(k\)-prime splits completely in \(\Omega_0\).
This proves (A.130.S4), including actual membership in the
image rather than only membership in its closure.

We also use its ramification-coordinate version:
prescribed full classes at \(T_0(K)\), and prescribed classes
modulo the unramified subgroup at finitely many other places,
are realized precisely when they annihilate
\(H^1(k_{T_0}/K,\mathbf Z/e)\).
Its pairing at a new unramified prime sends a valuation-one
Kummer class to evaluation of the dual character on Frobenius.
This is the same finite Poitou--Tate sequence and is the upper
exact row used on page~15 of the cited proof.

\emph{A scalar class of odd prime-power order.}
First take \(A=\mu_e\), where \(e\mid m\), viewed over \(K\).
Fix the usual compatible local Hilbert pairings with values
in \(\mu_e\); choose their order so that pairing an
unramified Kummer class with a valuation-one class evaluates
the former on Frobenius.
Their global product is one, and the unramified subgroup
is isotropic.  We write cohomology additively.

We construct classes \(z_i\in H^1(k_S/K,\mu_e)\) and primes
\(\mathfrak p_i\), with distinct underlying \(k\)-primes
outside \(T\), such that
\[
\begin{gathered}
 z_i\text{ has valuation }1\bmod e\text{ at }\mathfrak p_i,\\
 z_i\text{ is unramified at every other place};\\
 (z_i)_{\mathfrak p}=\tfrac12y_{\mathfrak p}
                                     \quad(\mathfrak p\in T(K));\\
 (z_j)_{\mathfrak p_i}=0,\qquad
 (z_j)_{\sigma\mathfrak p_i}=-(z_i)_{\sigma\mathfrak p_i},\\
 j>i,\qquad 1\ne\sigma\in\operatorname{Gal}(K/k).
\end{gathered}
 \tag{A.130.S5}
\]
The inverse of two exists modulo \(e\).
Every underlying prime is chosen to split completely in
\(\Omega\), so its different \(\operatorname{Gal}(K/k)\)-conjugates
are different places of \(K\).

Here is the precise induction.  Add all the earlier prime
orbits to \(T\), obtaining \(T_i\), and prescribe the last
two lines of (A.130.S5), with zero at each earlier distinguished
\(\mathfrak p_j\).  These are all unramified classes.
Their duality obstruction is a functional on the finite group
\(H^1(k_{T_i}/K,\mathbf Z/e)\).
Its restriction to the characters of \(\Omega/K\) is zero.
Indeed the earlier prime orbits split in \(\Omega\), while
at \(T(K)\) the functional is one half of that of the actual
global class \(y\); reciprocity for \(y\) has no contribution
outside \(T(K)\), because there \(y\) is unramified or the
prime splits in \(\Omega\).

Let \(F_i/K\) be the finite abelian extension cut out by
\(H^1(k_{T_i}/K,\mathbf Z/e)\).
The negative of the obstruction is an element of
\(\operatorname{Gal}(F_i/K)\) fixing \(F_i\cap\Omega\).
It therefore lifts to an element of
\(\operatorname{Gal}(F_i\Omega/K)\) acting identically on
\(\Omega\).  Chebotarev chooses a new underlying prime
splitting \(\Omega/k\) and a place \(\mathfrak p_{i+1}\)
above it with this Frobenius.  One may exclude all earlier
primes and any additional finite set.
Prescribe valuation one at that place and zero ramification
at its other conjugates.  Its Frobenius functional cancels
the obstruction.  The ramification-coordinate exact sequence
then supplies \(z_{i+1}\), proving (A.130.S5).

For each \(i\) record the finite vector
\[
 \lambda_i=\bigl(
 z_i(\operatorname{Frob}_{\sigma\mathfrak p_i})
                \bigr)_{1\ne\sigma\in\operatorname{Gal}(K/k)}
        \in\prod_{\sigma\ne1}\mu_e.
 \tag{A.130.S6}
\]
The displayed local restrictions are unramified, so these
values are defined.  Choose \(i<j\) with
\(\lambda_i=\lambda_j\), and put \(z=z_i+z_j\).
At \(T(K)\) this is \(y\).  At the conjugates of
\(\mathfrak p_i\), other than \(\mathfrak p_i\) itself,
the last line of (A.130.S5) gives zero directly.

For completeness the signs at the other distinguished
prime follow from two literal Hilbert products.
Fix \(\sigma\ne1\), and write
\(u=z_i(\operatorname{Frob}_{\sigma\mathfrak p_i})\).
In the product for \((z_i,\sigma z_i)\), only
\(\mathfrak p_i\) and \(\sigma\mathfrak p_i\) can contribute:
at \(T(K)\) both arguments are unramified, and elsewhere
both are unramified except at those two places.
The value at \(\sigma\mathfrak p_i\) is \(u\), so the value
at \(\mathfrak p_i\) is \(u^{-1}\).
By (A.130.S5) at \(\sigma^{-1}\mathfrak p_i\),
\[
       (\sigma z_j)_{\mathfrak p_i}
                         =-(\sigma z_i)_{\mathfrak p_i}.
 \tag{A.130.S7}
\]
The product for \((z_i,\sigma z_j)\) consequently has value
\(u\) at \(\mathfrak p_i\), and value \(u^{-1}\) at
\(\sigma\mathfrak p_j\).
The latter is exactly
\(z_i(\operatorname{Frob}_{\sigma\mathfrak p_j})\).
Equality of the vectors (A.130.S6) says that
\(z_j(\operatorname{Frob}_{\sigma\mathfrak p_j})=u\).
Their product is one.  Thus
\(z_{\sigma\mathfrak p_j}=0\) for every \(\sigma\ne1\).
At each distinguished prime the ramification coordinate
of \(z\) is still one modulo \(e\).
The scalar local homomorphism has cyclic target \(\mu_e\),
so it is cyclic.  This proves (A.130.S2) in the scalar case.
The argument used finiteness, Hilbert reciprocity, and the
inverse of two; it did not use that \(\mathbf Z/e\) is a field.

\emph{The bounded-character extension in a compositum.}
The following exact group fact supplies the module induction.
Let \(G_1\) be finite, let \(B\) be finite abelian killed by \(e\),
and let
\[
 G=G_1\times_I B,\qquad
       \pi_1:G\longrightarrow G_1,\quad\pi_2:G\longrightarrow B
 \tag{A.130.S8}
\]
be a fibre product for surjections onto a common quotient \(I\).
Every character \(\chi:G\to\mathbf Z/e\) has an expression
\[
                 \chi=\chi_1\circ\pi_1+\chi_2\circ\pi_2.
 \tag{A.130.S9}
\]
To prove this, identify \(\ker\pi_1\) with the subgroup
\(\ker(B\to I)\) of \(B\).
Its character \(\chi|_{\ker\pi_1}\) extends to a character
of \(B\) with values in \(\Q/\Z\).
This finite-character extension follows, for example, by
adjoining generators one at a time and taking the required
roots in the divisible group \(\Q/\Z\).
Since \(eB=0\), the image of every such extension is killed
by \(e\).  It is therefore a character \(\chi_2:B\to\mathbf Z/e\).
Subtracting its pullback from \(\chi\) makes the restriction
to \(\ker\pi_1\) zero; the difference descends to \(\chi_1\).
This proves (A.130.S9) without any linear-disjointness assertion.

\emph{Induction on the underlying cyclic summands.}
The construction of \(z\) and property (A.130.S2) concern
only the constant coefficient group over \(K\), its local
images, and the permutation of the places of \(K\).
We prove this construction for every finite abelian
\(p\)-group over \(K\), without prescribing an action of
\(\operatorname{Gal}(K/k)\) on that group.
The original arithmetic action is used only in the
corestriction at the end.
Decompose the underlying abelian group of \(A\) as a direct
sum of cyclic \(p\)-groups, ordered by nondecreasing exponent.
This is a decomposition over \(K\), where the action is trivial;
no summand is asserted to be stable under
\(\operatorname{Gal}(K/k)\).
Whenever a scalar construction uses a Galois action on
one cyclic summand, transport the cyclotomic action
through its chosen identification with the appropriate
\(\mu_{p^e}\).  Its conclusion about zero restrictions
at conjugate places does not depend on this temporary
extension of the trivial \(G_K\)-action.
Write the next induction step as
\(A=A'\oplus\mu_e\), using a chosen isomorphism of the last
cyclic group with \(\mu_e\).  The exponent of \(A'\) divides \(e\).
Apply the induction assertion over \(K\) to \(y'\) and obtain
\(z'\).  Let \(K'/K\) be the finite abelian extension cut out
by this homomorphism, and let \(\widetilde K/k\) be its
normal closure.  The group
\(\operatorname{Gal}(\widetilde K/K)\) embeds in a product
of the conjugate groups \(\operatorname{Gal}(K'/K)\).
It is therefore abelian of exponent dividing \(e\).
It is unramified at \(T(K)\), and any further ramification
lies above primes splitting in \(\Omega/k\).

Put
\[
 \Omega'=\Omega\widetilde K,\qquad
 T'=T\cup\operatorname{Ram}(\widetilde K/k),\qquad
 S'=\operatorname{cs}(\Omega'/k)\cup T'.
 \tag{A.130.S10}
\]
In particular \(S'\subset S\).
Prescribe a scalar local class \(\eta\) equal to \(y''\)
on \(T(K)\) and zero on \(T'(K)\setminus T(K)\).
Every \(\mathbf Z/e\)-valued character of
\(\operatorname{Gal}(\Omega'/K)\) is a sum of a character
from \(\Omega/K\) and one from \(\widetilde K/K\), by
(A.130.S8)--(A.130.S9).
The first character pairs to zero with \(\eta\) by global
reciprocity for the original global component \(y''\):
outside \(T(K)\), it vanishes at split primes and pairs
trivially with unramified classes elsewhere.
The second character is unramified at \(T(K)\), where
\(\eta\) is unramified, and \(\eta\) is zero on the added places.
It too pairs to zero.
The exact image statement (A.130.S4) therefore gives
\[
       \widetilde y\in H^1(k_{S'}/K,\mu_e),
                 \qquad \widetilde y|_{T'(K)}=\eta .
 \tag{A.130.S11}
\]

Apply the scalar construction to \(\Omega'/K/k\) and
\(\widetilde y\), obtaining \(z''\).
Set \(z=z'+z''\), viewed in the actual direct sum \(A\).
At \(T(K)\) it equals \(y\).
At a ramified place of \(z'\) outside \(T(K)\), the scalar
component is zero, because that place belongs to
\(T'(K)\setminus T(K)\).
At a new ramified place of \(z''\), its underlying prime
splits in \(\Omega'\), hence in the normal closure of \(K'\).
The entire component \(z'\), at that place and all its
\(\operatorname{Gal}(K/k)\)-conjugates, is zero.
Thus at every new ramified place only one of the two
components is present, and its nontrivial conjugate
restrictions vanish.  This proves (A.130.S2) inductively.
If every cyclic summand has order \(m\), its scalar
construction has inertia coordinate one modulo \(m\)
at both distinguished primes.  At a new prime belonging
to that summand, all earlier components vanish because
the prime splits their normal closure, and all later
components are prescribed to vanish there.  The inertia
therefore has order \(m\) and generates the full local
cyclic image.

Finally use the actual corestriction
\(H^1(k_S/K,A)\to H^1(k_S/k,A)\), with the originally
given \(\operatorname{Gal}(K/k)\)-action on \(A\).
At \(T\) its local Mackey formula gives the first equality
in (A.130.S3).
Outside \(T\), a ramified prime splits in \(\Omega\) and
hence in \(K\).  The same formula is a sum over its
distinct \(K\)-places.  Property (A.130.S2) leaves only
one nonzero ramified term, whose image is cyclic; its
conjugate by the actual coefficient action still has
cyclic image.  In the free case the coefficient action
is an automorphism, so it preserves the asserted order
\(m\) of inertia and the equality of the inertia and
decomposition images.  If no ramification remains, the class is
unramified and is killed by a finite unramified cyclic
extension.  This proves (A.130.S3) and completes the
odd-prime-power extension.



\begin{proof}[Proof of A.130]
The restricted-support Poitou--Tate comparison (A.122.6)
holds for every finite module, with its actual Cartier dual.
In particular it holds for \(C\), for \(C^\vee\), and after
base change to \(L\).  The finiteness of the fixed-set dual
group, stabilization of its decreasing kernels, and
Chebotarev in the finite cocycle field give precisely
\[
 \operatorname{loc}_{S_0}H^1_{S_0\cup{
m cs}(F)}(\Q,A)
       =\left(\operatorname{res}_{S_0}
                      H^1(\operatorname{Gal}(F/\Q),A^\vee)\right)^\perp
 \tag{A.130.6}
\]
for each finite module \(A\) split by \(F\), with the
coefficient and ramified places in \(S_0\).
The sum of local invariant pairings is the duality used
here.  The same proof permits deletion of finitely many
test primes.  See the exact maps in A.117 and A.122 and
\cite[I, Theorem 4.10(b) and Corollary 4.15]{BIR_MilneADT}.

Necessity in (A.130.4) follows from global reciprocity:
an alternating form vanishes on every cyclic subgroup,
and the other contributions outside \(S\) inflate from
the unramified procyclic quotient, whose degree-two
finite torsion cohomology is zero.
The coefficient comparison in A.129.5 identifies the
resulting Brauer classes with (A.130.3).

\emph{An actual full cocycle image at one new prime.}
Take \(\eta\) in the right side of (A.130.4) and choose
\(\gamma^{(0)}\in\mathcal H_r\) inducing it.
Its finite cocycle field is cut out by
\(g\mapsto(\gamma^{(0)}(g),\bar g)\in C\rtimes G\).
Choose \(q_0\) splitting \(EF_{\gamma^{(0)}}\), outside
all previously used places.
In the actual module
\(C=(x^d-1)R[C_N]\otimes\mu_{3^r}\), let
\(t_0=x^d-1\), with a fixed compatible root coordinate.
Its translates generate \(C\) over \(R\), and its
additive order is \(3^r\).
Apply (A.130.6) to add a global class which is zero
at \(S\) and all ramification of \(EF_{\gamma^{(0)}}\),
and is the pure tame class \((t_0,0)\) at \(q_0\).
These prescriptions annihilate every finite dual class
because \(q_0\) is completely split.
All the additional ramification of the added class
can be restricted to \({\rm cs}(EF_{\gamma^{(0)}})\).
Call the resulting class \(\gamma\), and put
\(\Omega=EF_\gamma\).

The homomorphism \(\gamma|_{G_E}\) contains \(t_0\)
in its image and its image is stable under \(G\).
It therefore has image all of \(C\), including the
full exponent \(3^r\).  The restriction maps give
\[
             F_\gamma\cap E=L,\qquad
       \operatorname{Gal}(\Omega/L)
                       =\operatorname{Gal}(E/L)\times C.
 \tag{A.130.7}
\]

\emph{The finite affine prescription with its exact dual.}
Let \(T\) contain \(S\), \(q_0\), and every prime
ramified in \(\Omega\), and let \(\mathcal R\) be the
primes where \(\gamma\) ramifies outside
\(S\cup\{q_0\}\).
Each such prime splits \(E\).
Write its tame and unramified coordinates as
\(\gamma_p=(t_p,c_p)\in C\oplus C\).
We seek a class over \(L\), ramified only above
\(T\cup{\rm cs}(\Omega)\), unramified above \(T\),
zero above \(S\), whose corestriction is
\((0,-c_p)\) at \(p\in\mathcal R\) and zero at \(q_0\).

For \(p\in\mathcal R\cup\{q_0\}\), index its
\(L\)-primes by \(\sigma\in G\) and transport back
by \(\sigma^{-1}\), including the coefficient action.
The local unramified Frobenius values \(a_{p,\sigma}\)
must satisfy the literal corestriction equations
\[
 \sum_{\sigma\in G}a_{p,\sigma}=-c_p
       \ (p\in\mathcal R),\qquad
                    \sum_{\sigma\in G}a_{q_0,\sigma}=0.
 \tag{A.130.8}
\]
They define a nonempty finite affine \(R\)-module.
Over \(L\), \(C\) and \(C^\vee\) are trivial, and
(A.130.7) identifies the finite dual group literally with
\[
 \operatorname{Hom}(\operatorname{Gal}(E/L),C^\vee)
                \oplus\operatorname{Hom}_{R}(C,C^\vee).
 \tag{A.130.9}
\]
The first summand is zero on the possible nonzero
prescriptions, since those rational primes split \(E\).
For \(f:C\to C^\vee\) in the second summand put
\[
 B_f(v,w)=\operatorname{ev}(v,f(w)),\qquad
 B_f^\sigma(v,w)=\chi_r(\sigma)^{-1}B_f(\sigma v,\sigma w).
 \tag{A.130.10}
\]
The full local dual obstruction, in compatible
\(\mu_{3^r}\) coordinates, is
\[
       -\sum_{p\in\mathcal R\cup\{q_0\}}
                   \sum_{\sigma\in G}B_f^\sigma(a_{p,\sigma},t_p).
 \tag{A.130.11}
\]
This is the same actual transported local pairing as
A.122.14, now at exponent \(3^r\); its minus sign is
the pairing of an unramified first class with a tame
second class.  Rational primes splitting \(\Omega\)
still detect the entire restriction kernel over \(L\):
one applies Chebotarev to the normal closure over
\(\Q\) of the finite dual cocycle field.
Thus (A.130.9)--(A.130.11) use the precise allowed
new prime set.

Let \(V_0\) be the variation module of (A.130.8).
An \(f\) annihilates every variation at \(q_0\)
exactly when
\[
           B_f^\sigma(v,t_0)=B_f(v,t_0)
                         \quad(v\in C,\ \sigma\in G).
 \tag{A.130.12}
\]
Comparing the equations for \(\sigma\tau\) and
\(\tau\) makes the same equality hold with
\(\tau t_0\) in the second argument.
The orbit of \(t_0\) generates over \(R\), so the
annihilator of all variations is exactly
\(\operatorname{Hom}_G(C\otimes_R C,\mu_{3^r})\).
A.129 computes this group: it consists of the
included mod-three alternating forms and nothing else.

For such a form \(B\), (A.130.11) is constant on
the affine prescription space, with value
\[
                         -\sum_{p\in\mathcal R}B(t_p,c_p).
 \tag{A.130.13}
\]
Global reciprocity and the assumed old sums make
this zero, since the value at \(q_0\) is pure tame
and the other unused contributions are unramified.
This gives an actual affine solution to all dual
conditions.  To spell out the integral finite
duality step, write \(D\) for the homomorphism from
local prescriptions to the Pontryagin dual of
\(\operatorname{Hom}_R(C,C^\vee)\) given by
(A.130.11).  Its variation image \(D(V_0)\) is a
finite subgroup.  The preceding calculation says
that \(D(a^{(0)})\) annihilates its annihilator
for any one affine solution \(a^{(0)}\).
Finite abelian double-annihilator duality implies
\(-D(a^{(0)})\in D(V_0)\).
This argument uses finite group duality and does
not require a vector-space complement over \(R\).
The corresponding variation gives the desired
class \(y\) over \(L\) by (A.130.6).

\emph{The odd-power cyclic correction and preservation of precision.}
Apply the preceding odd-power cohomological lemma
to \(\Omega/L/\Q\), exponent \(3^r\), module \(C\),
set \(T\), and the actual \(y\) just constructed.
The field \(L\) contains \(\mu_{3^r}\);
the class \(y\) is unramified at \(T\) and is zero
over every prime ramified in \(L/\Q\), over three,
and over the infinite place, all of which are in \(S\).
The lemma gives \(x\in H^1(\Q,C)\) with
\(x|_T=(\operatorname{cor}_{L/\Q}y)|_T\), cyclic
outside \(T\), and ramified only in
\(T\cup{\rm cs}(\Omega)\).
The underlying module \(C\) is free over \(R\), so the
stronger conclusion of the lemma applies: at a new ramified
prime its inertia has order \(3^r\) and generates its
entire local image.
Set \(\gamma'=\gamma+x\).
At every old place its class remains exactly
\(\eta_w\) modulo \(3^r\), because \(x_w=0\).
At \(p\in\mathcal R\) it is precisely \((t_p,0)\).
At a new ramified prime outside \(T\), the original
\(\gamma\) is locally zero because the prime
splits its cocycle field in \(\Omega\); hence
\(\gamma'_p=x_p\) has its local image generated by inertia.
All these new primes split \(E\).
At every earlier member of \(\mathcal R\), the class
\((t_p,0)\) has that same generation property, even when
\(t_p\) has smaller order.  It also holds at \(q_0\).
Thus \(\gamma'\in\mathcal H_{r,\rm in}\), proving both
image equalities in (A.130.4).

For the stated split old plans the linear conditions
in (A.130.6) vanish, because every finite dual class
is locally zero where a nonzero prescription occurs.
Pure tame plans have zero alternating quadratic
value.  This proves the specified corollary.

Finally, actual local Kummer prescriptions give an
actual global relative norm-one representative by
the torus comparison of A.117 at precision \(3^r\).
Here the other new places are split and the other
uncontrolled places are unramified, whose finite
residue torus has vanishing \(H^1\).
The relative cubic full-decomposition argument at
three gives the Hasse norm principle and makes
the resulting torus obstruction globally zero.
The torus Kummer exact sequence therefore yields
(A.130.5).  This last map retains an actual element
of the given field, while making no assertion
about compatible choices as \(r\) varies.
\end{proof}

The retained exact certificate
\path{code/relative_cyclic_precision_certificate.py}
uses integer Smith forms for the transported distribution
annihilators at five pairs \((n,r)\), including precision
twenty-seven.  It checks 368 full affine systems modulo
nine through finite abelian image orders, and 1,092 actual
bounded characters of the compositum fibre products.
Its output is retained beside the program.  The global
existence and the scalar Hilbert products are proved in
the odd-power lemma and the argument above.
